\section{Method}
\label{sec:method}

\subsection{Design}

The study is a within-subjects design with one factor, notification modality,
at four levels.

\begin{description}[itemsep=1pt,topsep=2pt,font=\normalfont\itshape]
  \item[Immediate.] A toast appears in the lower right corner the moment the
    build completes, with a result summary and a close control. It persists
    until dismissed. This is the behaviour of the editors our participants
    reported using daily.
  \item[Batched.] Notifications are held and delivered at the next pause in
    typing longer than three seconds, as a single toast that may summarise more
    than one build.
  \item[Ambient.] A gutter marker beside the edited line changes colour on
    completion, green or red. No motion, no sound, no dismissal control.
  \item[Silent.] No notification. The build status is available in the status
    bar if the participant looks at it.
\end{description}

Condition order was counterbalanced with a balanced Latin square, and task
assignment to condition was rotated independently, so no participant saw the
same task twice and no task appeared disproportionately in any condition.
Figure~\ref{fig:design} shows one participant's session structure.

\begin{figure}[t]
  \centering
  \begin{tikzpicture}[
      node distance=3.5mm and 4mm,
      blk/.style={draw,rounded corners=1.5pt,align=center,font=\scriptsize,
                  minimum height=6.5mm,minimum width=13mm},
      flow/.style={-{Stealth[length=1.6mm]},semithick},
    ]
    \node[blk] (intro) {Consent\\and intake};
    \node[blk,right=of intro] (train) {Training\\task};
    \node[blk,right=of train] (b1) {Block 1};
    \node[blk,below=6mm of b1] (b2) {Block 2};
    \node[blk,left=of b2] (b3) {Block 3};
    \node[blk,left=of b3] (b4) {Block 4};
    \node[blk,below=6mm of b4] (debrief) {Debrief\\interview};

    \draw[flow] (intro) -- (train);
    \draw[flow] (train) -- (b1);
    \draw[flow] (b1) -- (b2);
    \draw[flow] (b2) -- (b3);
    \draw[flow] (b3) -- (b4);
    \draw[flow] (b4) -- (debrief);

    \node[font=\scriptsize,right=2mm of b1,align=left] {one condition,\\one task};
    \node[font=\scriptsize,right=2mm of b2,align=left] {NASA-TLX\\after each block};

    \begin{scope}[on background layer]
      \node[draw,dashed,rounded corners=2pt,fit=(b1)(b2)(b3)(b4),
            inner sep=2.5mm,
            label={[font=\scriptsize]above:counterbalanced blocks}] {};
    \end{scope}
  \end{tikzpicture}
  \caption{Session structure. Each participant completed four blocks, one per
  notification condition, in an order drawn from a balanced Latin square.}
  \label{fig:design}
\end{figure}

\subsection{Participants}

We recruited 36 professional developers through two regional software meetups
and a university alumni list, screening for at least two years of paid
development experience and daily use of a compiled language. Sessions ran 95
minutes and participants received a voucher worth roughly one hour of median
local contracting rate. Four additional participants were recruited and
excluded before analysis: two because of eye-tracker calibration failure and
two because a task harness defect corrupted their keystroke logs.
Table~\ref{tab:demographics} describes the analysed sample.

\begin{table}[t]
  \centering
  \footnotesize
  \caption{Participant demographics, $n = 36$. Experience is years of paid
  professional development. Percentages may not sum to 100 because of
  rounding.}
  \label{tab:demographics}
  \begin{tabular}{llrr}
    \toprule
    Attribute & Level & $n$ & \% \\
    \midrule
    \multirow{4}{*}{Age}
      & 22 to 29 & 11 & 30.6 \\
      & 30 to 39 & 15 & 41.7 \\
      & 40 to 49 &  7 & 19.4 \\
      & 50 and above & 3 & 8.3 \\
    \midrule
    \multirow{3}{*}{Gender}
      & Woman & 14 & 38.9 \\
      & Man & 20 & 55.6 \\
      & Non-binary or undisclosed & 2 & 5.6 \\
    \midrule
    \multirow{3}{*}{Experience}
      & 2 to 5 years & 13 & 36.1 \\
      & 6 to 10 years & 14 & 38.9 \\
      & More than 10 years & 9 & 25.0 \\
    \midrule
    \multirow{4}{*}{Primary language}
      & Rust & 9 & 25.0 \\
      & TypeScript & 12 & 33.3 \\
      & Java or Kotlin & 10 & 27.8 \\
      & C or C++ & 5 & 13.9 \\
    \midrule
    \multirow{2}{*}{Team size}
      & Fewer than 6 & 21 & 58.3 \\
      & 6 or more & 15 & 41.7 \\
    \bottomrule
  \end{tabular}
\end{table}

\subsection{Tasks and apparatus}

Each task was a defect localisation and repair in a 4{,}200 line Rust service,
seeded with a failure that required editing two files. Tasks were matched for
size by a pilot with eight developers not included in the main sample, and
their mean completion times in the pilot fell within 11 percent of each other.

The environment was a purpose-built editor with the visual affordances of a
mainstream tool, instrumented to log every keystroke, caret move, scroll,
window focus change, and build event at millisecond resolution. A 120~Hz remote
eye tracker recorded gaze, calibrated at the start of each block and validated
between blocks. Builds were artificially held to 42~s, drawn from a truncated
normal with a 4~s standard deviation, matched to the median clean build time
reported by participants for their own projects.

\subsection{Measures}

\emph{Resumption lag} is the interval from the build notification event to the
first keystroke that modifies the primary task file, excluding navigation.
Where a participant never left the primary file, the lag is zero and the trial
contributes to the denominator. \emph{Completion time} is from task start to
the first passing test run. \emph{Error rate} is edits later reverted by the
same participant within the same task, per hundred edits. \emph{Mental demand}
is the corresponding NASA-TLX subscale, collected after each block on a
100-point scale.

\subsection{Analysis}

We fitted linear mixed models with condition as a fixed effect and participant
and task as crossed random intercepts, which handles the counterbalancing
without discarding the repeated measurements. Degrees of freedom use the
Satterthwaite approximation. Planned contrasts compare each condition against
\emph{Immediate}, with Holm correction across the three comparisons per
outcome. Confidence intervals are 95 percent profile intervals and effect sizes
are Cohen's \cohensd{} computed from the model's residual standard deviation.
The analysis plan was fixed before data collection.
